\chapter{High-Level Architecture Overview}

\section{What is QSOE?}

QSOE (Quick and Secure Operating Environment) is an operating environment
that combines the formally verified seL4 microkernel with a QNX
Neutrino-style programming model.  The result is a system that provides
the proven, elegant synchronous message-passing API of QNX on top of the
strongest isolation and capability-security guarantees available in any
production microkernel.

The target platform is RISC-V 64-bit, initially running on
\texttt{qemu-riscv-virt}.

\section{Design goals}

\begin{enumerate}
    \item \textbf{QNX-compatible API} --- provide the core QNX IPC
        primitives (\texttt{ChannelCreate}, \texttt{ConnectAttach},
        \texttt{MsgSend}/\texttt{Receive}/\texttt{Reply},
        \texttt{MsgSendPulse}, \texttt{ThreadCreate}, timers,
        synchronization) as a C library linked by user-space programs.
    \item \textbf{Formally verified foundation} --- leverage seL4's
        proven correctness and capability model rather than building a
        new kernel from scratch.
    \item \textbf{Minimal trusted computing base} --- only the seL4
        kernel runs in supervisor mode; everything else, including
        \texttt{taskman}, is a user-space process.
    \item \textbf{POSIX user-space via musl libc} --- standard C library
        support through musl, patched for seL4.
\end{enumerate}

\section{System layers}

Figure~\ref{fig:layers} shows the overall software stack.

\begin{figure}[ht]
\centering
\begin{tikzpicture}[
    layer/.style={
        draw, rounded corners=3pt, minimum width=12cm,
        minimum height=1.2cm, align=center, font=\small
    },
    label/.style={font=\footnotesize\itshape, text=gray},
]
    \node[layer, fill=blue!8]                          (kern) {seL4 microkernel};
    \node[layer, fill=green!8,  above=0.3cm of kern]   (boot) {elfloader / startup};
    \node[layer, fill=yellow!8, above=0.3cm of boot]   (rt)   {sel4runtime + musl libc};
    \node[layer, fill=orange!8, above=0.3cm of rt]     (tm)   {\texttt{taskman} --- process, memory, path, system services};
    \node[layer, fill=red!6,    above=0.3cm of tm]     (api)  {QNX-compatible C library (Channels, Connections, Messages)};
    \node[layer, fill=purple!6, above=0.3cm of api]    (app)  {User applications / resource managers};

    \node[label, left=0.3cm of kern]  {S-mode};
    \node[label, left=0.3cm of boot]  {M/S-mode};
    \node[label, left=0.3cm of rt]    {U-mode};
    \node[label, left=0.3cm of tm]    {U-mode};
    \node[label, left=0.3cm of api]   {U-mode};
    \node[label, left=0.3cm of app]   {U-mode};
\end{tikzpicture}
\caption{QSOE software stack}
\label{fig:layers}
\end{figure}

\section{Mapping QNX concepts to seL4}

The fundamental insight behind QSOE is that seL4 already provides the
low-level mechanisms that correspond closely to QNX's IPC model:

\begin{table}[ht]
\centering
\begin{tabular}{@{}ll@{}}
    \toprule
    \textbf{QNX concept}             & \textbf{seL4 foundation} \\
    \midrule
    Channel                          & seL4 Endpoint \\
    Connection (via \texttt{ConnectAttach}) & Capability to an Endpoint \\
    \texttt{MsgSend}/\texttt{MsgReceive}/\texttt{MsgReply} & seL4 synchronous IPC (\texttt{seL4\_Call}/\texttt{seL4\_Recv}/\texttt{seL4\_Reply}) \\
    Pulse (\texttt{MsgSendPulse})    & seL4 Notification \\
    Connection ID (coid)             & Capability badge \\
    Resource Manager                 & User-space server with capabilities \\
    \bottomrule
\end{tabular}
\caption{QNX--seL4 concept mapping}
\label{tab:mapping}
\end{table}
